\scanpage{33}
Soit $f:X\longrightarrow Y$ loc. de t.f.

On définit $D_{\mathrm{parf}}(f)$ comme la catégorie \inserted{triangulée} des complexes de $\mathcal O_X$-Modules \uline{parfaits} rel. à $f$, et on définit $K^\bullet_{\mathrm{parf}}(f)$ en conséquence.

Soit $Z\subset X$ un fermé d'un préschéma.

On définit $D_{\mathrm{parf},Z}(X)$ comme la catégorie des complexes de $\mathcal O_X$-modules qui sont parfaits et à cohomologie à support $\subset Z$. D'où aussi $K_Z^\bullet(X)$.
On définit de $=$ $D_{\mathrm{parf},\Phi}(X)$, $K_\Phi^\bullet(X)$, $\Phi$ une ``famille $\Phi$'', \unsure{d'où}
\[
\left\{\begin{array}{l}
\deleted{\unsure{K}}D_{\mathrm{parf},\Phi}(X)=\displaystyle\mathop{\mathrm{Lim}}\limits_{\longrightarrow\atop Z\in\Phi}D_{\mathrm{parf},Z}(X)\\
K_\Phi^\bullet(X)=\displaystyle\mathop{\underline{\lim}}\limits_{Z\in\Phi}K^\bullet_{\unsure{\Phi|Z}}(X),
\end{array}\right.
\]

Soit $i:Z\mathbin{\text{---}}X$ une imm. fermée. On a un foncteur can
\[
\mathrm Ri_*:D_{\mathrm{parf}}(\deleted{\unsure{f}}\inserted{i})\longrightarrow D_Z(X)\qquad\text{d'où}
\]
\[
i_!:K^\bullet(\deleted{\unsure{f}}\inserted{i})\longrightarrow K_Z^\bullet(X)
\]
Soit $\deleted{\unsure{f_n}}i_n:Z_n\longrightarrow X$ le $n$-ième vois. inf. de $Z$ dans $X$, on a alors
\[
\mathop{\mathrm{Lim}}\limits_{\longrightarrow\atop n}D_{\mathrm{parf}}(\deleted{\unsure{f_n}}\inserted{i_n})\longrightarrow D_Z(X)
\]
\uline{le foncteur est pl. fidèle.}
\margin{32}
